\section*{अध्याय \ref{ch:g11:minerals} --- खनिज}

\begin{solution}{exo:g11:minerals:1}
\omterm{def:g11:minerals:cleavage}{विदलन} किसी खनिज की समतल तलों के अनुदिश टूटने की प्रवृत्ति है, जो उन
दिशाओं के समांतर होते हैं जिनमें उसके परमाणुओं के बीच के बंध कमज़ोर हैं; \omterm{def:g11:minerals:cleavage}{विभंग}
अनियमित सतहों के अनुदिश टूटना है, ऐसे खनिज में जिसमें कोई कमज़ोर तल नहीं।
क्वार्ट्ज़ में कोई \omterm{def:g11:minerals:cleavage}{विदलन} नहीं: वह शंखाभ \omterm{def:g11:minerals:cleavage}{विभंग} दिखाता है, टूटे काँच जैसी चिकनी
वक्र सतहों के साथ। अभ्रक का एक \omterm{def:g11:minerals:cleavage}{विदलन} होता है, पूर्ण: वह परतों में
उतरता है।
\end{solution}

\begin{solution}{exo:g11:minerals:2}
किसी दूसरे तत्त्व का अंश मात्र क्वार्ट्ज़ का रंग बदलने के लिए काफ़ी है, जो
पारदर्शी, सफ़ेद, काला, गुलाबी, नीला या बैंगनी हो सकता है; ओलिवीन लगभग हमेशा
हरा होता है। अधिक भरोसेमंद: \omterm{def:g11:minerals:cleavage}{विदलन} (या \omterm{def:g11:minerals:cleavage}{विभंग}), कठोरता, \omterm{def:g11:minerals:lustre}{द्युति},
खुलकर बढ़े क्रिस्टलों का \omterm{def:g11:minerals:habit}{स्वभाव}, घनत्व।
\end{solution}

\begin{solution}{exo:g11:minerals:3}
(1) ओलिवीन। (2) बायोटाइट। (3) ऑर्थोक्लेज़ (एक फ़ेल्डस्पार: दो \omterm{def:g11:minerals:cleavage}{विदलन}, कठोरता
6)। (4) क्वार्ट्ज़ (कोई \omterm{def:g11:minerals:cleavage}{विदलन} नहीं, कठोरता 7)।
\end{solution}

\begin{solution}{exo:g11:minerals:4}
(1) पाँच परमाणु: केंद्र पर एक सिलिकन और एक चतुष्फलक के कोनों पर चार ऑक्सीजन।
(2) $+4 + 4 \times (-2) = -4$: समूह \ce{SiO4^4-} है।
\end{solution}

\begin{solution}{exo:g11:minerals:5}
ओलिवीन: पृथक चतुष्फलक। पाइरॉक्सीन: एकल शृंखलाएँ। ऐम्फ़िबोल: द्वि-शृंखलाएँ।
अभ्रक: परतें। फ़ेल्डस्पार और क्वार्ट्ज़: ढाँचे।
\end{solution}

\begin{solution}{exo:g11:minerals:6}
(1) क्वार्ट्ज़ का मोलर द्रव्यमान $M(\ce{SiO2}) = 28.085 + 2 \times 15.999 = \qty{60.083}{g/mol}$;
ऑक्सीजन: $31.998/60.083 \approx \qty{53.3}{\percent}$।
(2) फ़ॉर्स्टेराइट का मोलर द्रव्यमान $M(\ce{Mg2SiO4}) = 2 \times 24.305 + 28.085 + 4 \times 15.999 =
\qty{140.691}{g/mol}$; ऑक्सीजन: $63.996/140.691 \approx \qty{45.5}{\percent}$।
(3) भूपर्पटी, \qty{46}{\percent} पर, उसी कोटि की है: वह लगभग पूरी
ऑक्सीजन-समृद्ध खनिजों, \omterm{def:g11:minerals:silicate}{सिलिकेटों}, से बनी है।
\end{solution}

\begin{solution}{exo:g11:minerals:7}
सफ़ेद कण के दो \omterm{def:g11:minerals:cleavage}{विदलन} हैं और कठोरता 5.5 और 6.5 के बीच: एक
फ़ेल्डस्पार (6)। धूसर कण 6.5 से अधिक कठोर है और बिना किसी \omterm{def:g11:minerals:cleavage}{विदलन} के एक वक्र,
शंखाभ \omterm{def:g11:minerals:cleavage}{विभंग} दिखाता है: क्वार्ट्ज़ (7)।
\end{solution}

\begin{solution}{exo:g11:minerals:8}
(1) प्रति सिलिकन $4 - 2/2 = 3$ ऑक्सीजन: \ce{SiO3}। (2) \ce{(Si,Al)2O6} में
छह ऑक्सीजनों के लिए दो सिलिकन (या ऐलुमिनियम) परमाणु हैं, प्रति सिलिकन तीन: यह
दो बार \ce{SiO3} है।
\end{solution}

\begin{solution}{exo:g11:minerals:9}
$\rho = 0.55 \times 2.69 + 0.35 \times 3.4 + 0.10 \times 3.32 = 1.480 + 1.190
+ 0.332 \approx \qty{3.00}{g/cm^3}$। घन का आयतन
\qty{1000}{cm^3} है, इसलिए द्रव्यमान लगभग \qty{3.00}{kg}।
\end{solution}

\begin{solution}{exo:g11:minerals:10}
(1) $0.75 \times 3.32 + 0.25 \times 3.4 = 2.49 + 0.85 = \qty{3.34}{g/cm^3}$।
(2) $(3.34 - 3.00)/3.00 \approx 0.11$: लगभग \qty{11}{\percent} अधिक घना।
(3) \omterm{def:g11:ocean-ridges:oceanic-lithosphere}{महासागरीय भूपर्पटी} के बेसाल्ट और \omterm{def:g11:ocean-ridges:gabbro}{गैब्रो} प्रावार के \omterm{def:g11:ocean-ridges:peridotite}{पेरिडोटाइट} से
हल्के हैं: वे उसके ऊपर टिके रहते हैं।
\end{solution}

\begin{solution}{exo:g11:minerals:11}
(1) प्लेजियोक्लेज़ $168/400 = 0.42$, पाइरॉक्सीन $0.33$, ओलिवीन $0.15$, छिद्र
$0.10$। (2) प्रस्ताव बिना रंध्रों की चट्टान के लिए लागू होता है; यहाँ आयतन का
दसवाँ भाग खाली है। (3) ठोस भाग में अंश $0.42/0.90
\approx 0.467$, $0.367$ और $0.167$ हैं:
$\rho_{\text{ठोस}} \approx 0.467 \times 2.69 + 0.367 \times 3.4 + 0.167
\times 3.32 \approx \qty{3.06}{g/cm^3}$। अपने छिद्रों सहित, जो आयतन जोड़ते हैं
पर कोई द्रव्यमान नहीं, चट्टान का घनत्व $0.42 \times 2.69 + 0.33 \times 3.4 +
0.15 \times 3.32 \approx \qty{2.75}{g/cm^3}$ है, यानी $0.9 \times 3.06$।
\end{solution}

\begin{solution}{exo:g11:minerals:12}
(1) ऑक्सीजन $464/16.0 = 29.0$~mol; सिलिकन $283/28.1 \approx 10.1$; ऐलुमिनियम
$84/27.0 \approx 3.1$; लोहा $52/55.8 \approx 0.93$; कैल्सियम $46/40.1 \approx
1.15$; मैग्नीशियम $28/24.3 \approx 1.15$; सोडियम $23/23.0 = 1.0$; पोटैशियम
$15/39.1 \approx 0.38$। (2) कुल मिलाकर लगभग $46.8$~mol परमाणु, जिनमें से
$29.0/46.8 \approx 0.62$: लगभग तीन में से दो परमाणु ऑक्सीजन हैं।
(3) बड़े ऑक्सीजन आयन अधिकांश परमाणु बनाते हैं, और आयतन का उससे भी
अधिक भाग: परमाणुओं के पैमाने पर भूपर्पटी ऑक्सीजन आयनों का एक संकुलन है, जिसमें
छोटे सिलिकन और धातु आयन उनके बीच की जगहों में हैं।
\end{solution}

\begin{solution}{exo:g11:minerals:13}
(1) साझा कोनों की माध्य संख्या $(2 + 3)/2 = 2.5$ है, इसलिए प्रति सिलिकन
$4 - 2.5/2 = 2.75$ ऑक्सीजन हैं। (2) ग्यारह ऑक्सीजनों के लिए चार सिलिकन:
\ce{Si4O11}। (3) \ce{Si7AlO22} में 22 ऑक्सीजनों के लिए आठ चतुष्फलक (सात सिलिकन, एक
ऐलुमिनियम) हैं, यानी प्रति चतुष्फलक $22/8 = 2.75$, या
दो बार \ce{Si4O11}। दोनों \ce{OH} समूह चतुष्फलकों का भाग नहीं हैं।
\end{solution}

\begin{solution}{exo:g11:minerals:14}
$0.20 \times 2.62 + 0.10 \times 2.56 + 0.50 \times 2.65 + 0.10 \times 3.09 +
0.10 \times 3.23 = 0.524 + 0.256 + 1.325 + 0.309 + 0.323 \approx
\qty{2.74}{g/cm^3}$, \qty{2.7}{g/cm^3} से ऊपर; उदाहरण का ग्रेनाइट,
क्वार्ट्ज़ और पोटैशियम फ़ेल्डस्पार से अधिक समृद्ध, लगभग \qty{2.59}{g/cm^3} का है,
\qty{2.6}{g/cm^3} से नीचे। पश्चिम की चट्टानें अधिक घनी हैं क्योंकि उनमें अधिक
प्लेजियोक्लेज़ और अधिक बायोटाइट और हॉर्नब्लेंड है, जो लोहे और मैग्नीशियम से समृद्ध हैं;
उच्च-सिलिका ग्रेनाइटों में अधिक क्वार्ट्ज़ और पोटैशियम फ़ेल्डस्पार है, जो सबसे हल्के
शैलकारी खनिज हैं।
\end{solution}

\begin{solution}{exo:g11:minerals:15}
(1) ऊपरी भूपर्पटी: $66.6 \times 0.467 \approx \qty{31.1}{\percent}$ सिलिकन;
समग्र भूपर्पटी: $60.6 \times 0.467 \approx \qty{28.3}{\percent}$। (2) ऊपरी
भूपर्पटी में, जो सिलिका से अधिक समृद्ध है, अधिक क्वार्ट्ज़ और फ़ेल्डस्पार होने चाहिए, और
लोहे और मैग्नीशियम से समृद्ध खनिज कम; गहरी भूपर्पटी में इसका उलटा। (3) क्वार्ट्ज़
और फ़ेल्डस्पार सबसे हल्के शैलकारी खनिज हैं: ऊपरी भूपर्पटी गहरी भूपर्पटी से
हल्की होनी चाहिए।
\end{solution}

\begin{solution}{pb:g11:minerals:1}
\textbf{1.} क्वार्ट्ज़: कोई \omterm{def:g11:minerals:cleavage}{विदलन} नहीं, शंखाभ \omterm{def:g11:minerals:cleavage}{विभंग}, काचाभ \omterm{def:g11:minerals:lustre}{द्युति}, कील से
अधिक कठोर (7)।

\textbf{2.} ऑर्थोक्लेज़, एक पोटैशियम फ़ेल्डस्पार: दो \omterm{def:g11:minerals:cleavage}{विदलन}, मोह्स पैमाने पर
कठोरता 6।

\textbf{3.} प्लेजियोक्लेज़ (यहाँ ऑलिगोक्लेज़), सोडियम और
कैल्सियम वाला फ़ेल्डस्पार।

\textbf{4.} D बायोटाइट है, E मस्कोवाइट: दो अभ्रक। उनके चतुष्फलक
परतें बनाते हैं, और परतों के बीच के बंध कमज़ोर हैं: वे इसी एक
पूर्ण \omterm{def:g11:minerals:cleavage}{विदलन} के अनुदिश फटते हैं।

\textbf{5.} क्वार्ट्ज़ और दोनों फ़ेल्डस्पार (A, B, C) ढाँचे हैं; बायोटाइट
(D) में लोहा और मैग्नीशियम है।

\textbf{6.} क्वार्ट्ज़ $150/500 = 0.30$; ऑर्थोक्लेज़ $0.35$; ऑलिगोक्लेज़ $0.25$;
बायोटाइट $0.06$; मस्कोवाइट $0.04$।

\textbf{7.} जब कण बिना किसी क्रम के बिखरे हों, तो किसी खनिज पर पड़ने वाले बिंदुओं का अंश
फलक पर उसके द्वारा ढके क्षेत्रफल के अंश का अनुमान देता है, और वह
उसके द्वारा भरे आयतन के अंश के बराबर है।

\textbf{8.} $0.35 + 0.25 = 0.60$: \qty{60}{\percent} फ़ेल्डस्पार।

\textbf{9.} $1/500 = 0.002$ से, यानी \qty{0.2}{\percent}।

\textbf{10.} हाँ: फ़ेल्डस्पार, क्वार्ट्ज़ और अभ्रक।

\textbf{11.} $\rho = \phi_1 \rho_1 + \phi_2 \rho_2 + \dots + \phi_k \rho_k$,
जहाँ $\phi_1 + \dots + \phi_k = 1$।

\textbf{12.} क्वार्ट्ज़ $0.30 \times 2.62 = 0.786$; ऑर्थोक्लेज़ $0.35 \times 2.56 =
0.896$; ऑलिगोक्लेज़ $0.25 \times 2.65 = 0.6625$; बायोटाइट $0.06 \times 3.09
\approx 0.185$; मस्कोवाइट $0.04 \times 2.82 \approx 0.113$ (सभी
\unit{g/cm^3} में)।

\textbf{13.} $0.786 + 0.896 + 0.6625 + 0.185 + 0.113 \approx
\qty{2.64}{g/cm^3}$।

\textbf{14.} ऑर्थोक्लेज़, $0.896$ पर, क्योंकि वह सबसे प्रचुर है। बायोटाइट
और मस्कोवाइट, \qty{3.09}{} और \qty{2.82}{g/cm^3} पर, दूसरे खनिजों से
अधिक घने हैं और चट्टान का घनत्व अकेले क्वार्ट्ज़ और फ़ेल्डस्पार के घनत्व से
ऊपर उठाते हैं।

\textbf{15.} \qty{2640}{kg/m^3}: ग्रेनाइट के \qty{1}{m^3} का द्रव्यमान
लगभग \qty{2.64}{t} है।

\textbf{16.} $100 \times 2.643 \approx \qty{264.3}{g}$।

\textbf{17.} खंड \qty{100}{cm^3} पानी हटाता है, जिसका द्रव्यमान
\qty{100}{g} है: तुला $264.3 - 100 = \qty{164.3}{g}$ पढ़ती है।

\textbf{18.} $264.3/(264.3 - 164.3) = 264.3/100 \approx \qty{2.64}{g/cm^3}$।

\textbf{19.} दोनों के बीच: \qty{2.69}{g/cm^3} के औसत से हल्का,
सिलिका और पोटैशियम से सबसे समृद्ध ग्रेनाइटों के \qty{2.6}{g/cm^3} से ऊपर।
अधिक प्लेजियोक्लेज़ और अधिक बायोटाइट या हॉर्नब्लेंड, और कम क्वार्ट्ज़ और
पोटैशियम फ़ेल्डस्पार, किसी ग्रेनाइट को अधिक घना बनाएँगे।

\textbf{20.} ग्रेनाइट का घनत्व लगभग
\textbf{\qty{2.64}{g/cm^3}} है।
\end{solution}
